\section{Method}

\subsection{Problem Setting}

We target aerial small-object detection under a fixed RetinaNet + FPN detector. The backbone is the only algorithmic variable in the main experiments.

\subsection{Frequency-Decoupled Hybrid Mamba}

Given an intermediate feature map $x$, we compute a lightweight low-pass component and define the residual high-frequency component as the complement. The low-frequency part is routed to the Mamba branch, while the high-frequency part is processed by a lightweight detail branch.

\subsection{Detail Branch}

The detail branch uses parallel depthwise $3 \times 3$ and $5 \times 5$ convolutions, followed by an edge-aware gate and pointwise fusion. This branch is deliberately small: its role is to restore local high-frequency evidence rather than to replace the global modeling capacity of the Mamba branch.

\subsection{Backbone Variants}

We evaluate three backbone variants under the same detector:

\begin{itemize}[leftmargin=1.5em]
  \item TinyViM baseline
  \item HybridMamba-Base, which keeps the low-frequency-only Mamba path
  \item HybridMambaDet, which adds the lightweight high-frequency detail branch
\end{itemize}

